\documentclass{beamer}

% 引入必要的宏包
\usepackage[utf8]{inputenc}
\usepackage{ctex} % 支持中文排版
\usepackage{amsmath}

% 设置 Beamer 主题和配色
\usetheme{Madrid}
\usecolortheme{default}

% 幻灯片首页信息
\title{决策树是非线性模型的数学证明}
\subtitle{基于线性算子定义的严格推导}
\author{演示者}
\date{\today}

\begin{document}

% 封面页
\frame{\titlepage}

% 第1页：线性的严格定义
\begin{frame}
\frametitle{1. 线性的严谨数学定义}
一个模型或函数 $f(x)$ 是线性的，当且仅当对于任意标量 $\alpha$ 和任意输入特征向量 $x, y$，严格满足以下两个条件：
\vspace{0.5cm}
\begin{itemize}
    \item \textbf{齐次性 (Homogeneity):} 
    $$f(\alpha x) = \alpha f(x)$$
    \item \textbf{可加性 (Additivity):} 
    $$f(x + y) = f(x) + f(y)$$
\end{itemize}
\vspace{0.5cm}
\begin{block}{证明思路}
采用\textbf{反证法}：只要构造出一组特定输入，证明决策树违反上述任意一个条件，即可证明其为非线性模型。
\end{block}
\end{frame}

% 第2页：决策树的数学表达
\begin{frame}
\frametitle{2. 决策树的数学表达式}
决策树本质上是特征空间上的\textbf{分段常数函数}（阶跃函数的组合）。其预测函数形式化表示为：
$$f(x) = \sum_{m=1}^{M} c_m \mathbb{I}(x \in R_m)$$
\vspace{0.2cm}
\textbf{其中：}
\begin{itemize}
    \item $M$: 叶子节点的总数量
    \item $R_m$: 第 $m$ 个叶子节点对应的特征空间划分区域
    \item $c_m$: 落在该区域内的样本的常数预测值
    \item $\mathbb{I}(\cdot)$: 指示函数（条件成立为 1，否则为 0）
\end{itemize}
\end{frame}

% 第3页：证明齐次性失效
\begin{frame}
\frametitle{3. 构造性证明：破坏齐次性}
假设有一棵单层回归决策树模型：
\begin{equation*}
f(x) =
\begin{cases}
  10 & \text{if } x \leq 5 \\
  20 & \text{if } x > 5
\end{cases}
\end{equation*}

我们来验证齐次性 $f(\alpha x) = \alpha f(x)$。令输入 $x = 2$，标量 $\alpha = 3$：
\vspace{0.2cm}
\begin{itemize}
    \item \textbf{等式左侧：} $f(\alpha x) = f(3 \times 2) = f(6)$。因为 $6 > 5$，所以 \alert{$f(6) = 20$}。
    \item \textbf{等式右侧：} $\alpha f(x) = 3 \times f(2)$。因为 $2 \leq 5$，所以 $f(2) = 10$，结果为 $3 \times 10 = \alert{30}$。
\end{itemize}
\vspace{0.3cm}
显然 $20 \neq 30$，\textbf{齐次性假设不成立}。
\end{frame}

% 第4页：证明可加性失效
\begin{frame}
\frametitle{4. 构造性证明：破坏可加性}
沿用上述同一棵决策树：
\begin{equation*}
f(x) =
\begin{cases}
  10 & \text{if } x \leq 5 \\
  20 & \text{if } x > 5
\end{cases}
\end{equation*}

我们来验证可加性 $f(x + y) = f(x) + f(y)$。令输入 $x = 2, y = 2$：
\vspace{0.2cm}
\begin{itemize}
    \item \textbf{等式左侧：} $f(x + y) = f(2 + 2) = f(4)$。因为 $4 \leq 5$，所以 \alert{$f(4) = 10$}。
    \item \textbf{等式右侧：} $f(x) + f(y) = f(2) + f(2) = 10 + 10 = \alert{20}$。
\end{itemize}
\vspace{0.3cm}
显然 $10 \neq 20$，\textbf{可加性假设同样不成立}。
\end{frame}

% 第5页：总结
\begin{frame}
\frametitle{5. 最终结论}
\begin{block}{核心数学结论}
由于决策树模型的预测算子既不满足齐次性，也不满足可加性，因此在严格的数学定义下，\textbf{决策树是一种非线性模型}。
\end{block}
\vspace{0.5cm}
\begin{block}{物理与几何意义}
公式中的指示函数 $\mathbb{I}(\cdot)$ 在分裂点处会导致一阶导数不存在（发生跳变）。这种不连续的“阶跃”特质，正是决策树能够无需手动构造多项式，即可天然拟合高阶交互特征和非线性边界的根本原因。
\end{block}
\end{frame}

\end{document}